\section{More Facts About Singular Values}

Recall the maximum- and minimum-gain definitions
\[
\sigma_{\max}(X)=\max_{\|z\|=1}\|Xz\|,
\qquad
\sigma_{\min}(X)=\min_{\|z\|=1}\|Xz\|
\]

\textbf{(a)}
\textbf{True}.

Let \(e_i\) be the \(i\)th standard basis vector. The norm of
row \(i\) of \(X\) is
\[
\left(\sum_{j=1}^n |X_{ij}|^2\right)^{1/2}
=\|X^Te_i\|
\]
Since \(\|e_i\|=1\) and \(X\) and \(X^T\) have the same singular values,
\[
\|X^Te_i\|\leq\sigma_{\max}(X^T)=\sigma_{\max}(X)
\]
Taking the maximum over \(i\) proves the statement.

\textbf{(b)}
\textbf{False}.

If \(X\in\mathbb{R}^{n\times n}\) is rank deficient, then
some unit vector \(z\) satisfies \(Xz=0\), so
\[
\sigma_{\min}(X)=0
\]
This can occur even when every row has positive norm. For example,
\[
X=
\begin{bmatrix}
1&0\\
1&0
\end{bmatrix}
\]
Both row norms equal \(1\), but \(X\) has rank one, so the claimed
inequality would give \(0\geq1\).

\textbf{(c)}
\textbf{True}.

We know that
\[
\|XY\|
\leq\|X\|\|Y\|
\]
So for any \(z\),
\[
\|XYz\|
\leq\sigma_{\max}(X)\|Yz\|
\leq\sigma_{\max}(X)\sigma_{\max}(Y)\|z\|
\]
Maximizing over all unit vectors \(z\) gives
\[
\boxed{\sigma_{\max}(XY)
\leq\sigma_{\max}(X)\sigma_{\max}(Y)}
\]

\textbf{(d)}
\textbf{True}.

For any \(z\),
\[
\|XYz\|
\geq\sigma_{\min}(X)\|Yz\|
\geq\sigma_{\min}(X)\sigma_{\min}(Y)\|z\|
\]
Then, minimizing over all unit vectors \(z\) gives
\[
\boxed{\sigma_{\min}(XY)
\geq\sigma_{\min}(X)\sigma_{\min}(Y)}
\]

\textbf{(e)}
\textbf{True}.

For every unit vector \(z\), (reverse) triangle inequality
gives
\[
\|(X+Y)z\|
\;=\|Xz+Yz\|
\geq\|Xz\|-\|Yz\|
\geq\sigma_{\min}(X)-\sigma_{\max}(Y)
\]

So taking the minimum over \(z\) gives
\[
\boxed{\sigma_{\min}(X+Y)
\geq\sigma_{\min}(X)-\sigma_{\max}(Y)}
\]
